% 第6章 6.6.2：基于谱分解的滤波在单位权三点路径上的逐模态执行。
% f=(3,1,2)，g(lambda)=1/(1+lambda)，特征对 (0,(1,1,1)/sqrt3)、(1,(1,0,-1)/sqrt2)、(3,(1,-2,1)/sqrt6)。
% 输入成分 fhat_k u_k：(2,2,2)、(1/2,0,-1/2)、(1/2,-1,1/2)，相加为 f=(3,1,2)。
% 输出成分 yhat_k u_k=g(lambda_k) fhat_k u_k：(2,2,2)、(1/4,0,-1/4)、(1/8,-1/4,1/8)，相加为 y=(19/8,7/4,15/8)。
% 柱高 4mm/单位；蓝色为输入成分，绿色为输出成分。
\begin{tikzpicture}[x=1mm,y=1mm,font=\figurefont\footnotesize,text=FigureInk,
  base/.style={draw=FigureStroke,line width=.6pt},
  inbar/.style={draw=FigureBlue,fill=FigureBlueFill,line width=.7pt},
  outbar/.style={draw=FigureGreen,fill=FigureGreenFill,line width=.7pt},
  val/.style={font=\figurefont\fontsize{7}{8}\selectfont,text=FigureInk,inner sep=.6pt},
  note/.style={font=\FigureNoteFont\fontsize{7.5}{9}\selectfont,text=FigureMuted},
  arr/.style={-{Stealth[length=1.6mm]},draw=FigureStroke,line width=.7pt}]
\def\s{4}
% 一行柱状图：#1 x 起点，#2 基线 y，#3 样式，#4/#5/#6 三个值（数值），#7/#8/#9 三个标签
\newcommand{\bars}[9]{%
  \draw[base] (#1-0.5,#2)--(#1+20.5,#2);
  \foreach \i/\v/\t in {0/#4/{#7},1/#5/{#8},2/#6/{#9}}{%
    \pgfmathsetmacro{\xx}{#1+3+7*\i}%
    \pgfmathsetmacro{\hh}{\s*\v}%
    \ifdim \hh pt=0pt
      \node[val,anchor=south] at (\xx,#2) {$\t$};
    \else
      \draw[#3] (\xx-1.8,#2) rectangle (\xx+1.8,#2+\hh);
      \ifdim \hh pt>0pt
        \node[val,anchor=south] at (\xx,#2+\hh) {$\t$};
      \else
        \node[val,anchor=north] at (\xx,#2+\hh) {$\t$};
      \fi
    \fi}}
% 表头
\node[anchor=south] at (26,62) {输入的成分$\widehat f_k\bm u_k$};
\node[anchor=south] at (69,62) {输出的成分$\widehat y_k\bm u_k$};
% 第1行：k=1
\node[anchor=west] at (0,53) {$\lambda_1=0$};
\bars{16}{50}{inbar}{2}{2}{2}{2}{2}{2}
\draw[arr] (38.5,53)--node[above,val]{$\times g(0)=1$}(56.5,53);
\bars{59}{50}{outbar}{2}{2}{2}{2}{2}{2}
% 第2行：k=2
\node[anchor=west] at (0,42) {$\lambda_2=1$};
\bars{16}{42}{inbar}{0.5}{0}{-0.5}{\tfrac12}{0}{-\tfrac12}
\draw[arr] (38.5,42)--node[above,val]{$\times g(1)=\tfrac12$}(56.5,42);
\bars{59}{42}{outbar}{0.25}{0}{-0.25}{\tfrac14}{0}{-\tfrac14}
% 第3行：k=3
\node[anchor=west] at (0,28.5) {$\lambda_3=3$};
\bars{16}{28.5}{inbar}{0.5}{-1}{0.5}{\tfrac12}{-1}{\tfrac12}
\draw[arr] (38.5,28.5)--node[above,val]{$\times g(3)=\tfrac14$}(56.5,28.5);
\bars{59}{28.5}{outbar}{0.125}{-0.25}{0.125}{\tfrac18}{-\tfrac14}{\tfrac18}
% 分隔线与求和
\draw[draw=FigureGrid,line width=.6pt] (0,19.5)--(80,19.5);
\node[anchor=west,note] at (0,7) {按列相加};
\bars{16}{2}{inbar}{3}{1}{2}{3}{1}{2}
\bars{59}{2}{outbar}{2.375}{1.75}{1.875}{\tfrac{19}8}{\tfrac74}{\tfrac{15}8}
\node at (47.5,8) {$\bm f\ \longmapsto\ \bm y$};
\foreach \x in {16,59}{\foreach \i in {1,2,3}{\node[note,anchor=north] at (\x+3+7*\i-7,1.2) {$\i$};}}
\node[note,anchor=north] at (47.5,1.2) {顶点编号};
\end{tikzpicture}
